\section{Discussion \& Conclusion}
    본 장에서는 6장의 시뮬레이션 및 실험 결과를 종합하여 필터별 성능을 비교 분석하고, 연구의 한계 및 향후 과제를 제시한다.\\
    
    \subsection{Filter Performance Comparison}
        \subsubsection{Accuracy}
            시뮬레이션 및 실험 결과를 종합하면, 자세 추정 정확도는 일반적으로 EKF $>$ Madgwick $\approx$ Mahony $>$ CF 순으로 나타난다.\\
 
            EKF는 센서 잡음의 확률적 특성을 Kalman Gain에 반영하므로, 동적 환경에서도 최적에 가까운 추정 성능을 제공한다. 특히 자이로 바이어스를 온라인으로 추정함으로써 장시간 운용 시 드리프트 억제 효과가 두드러진다.\\
 
            Mahony Filter는 PI 바이어스 추정 덕분에 CF보다 드리프트에 강하며, Madgwick Filter는 단일 파라미터 $\beta$로 튜닝이 간편하다. CF는 구조가 단순하여 계산 비용이 가장 낮지만, 고정된 $\alpha$로 인해 동적 환경에서 성능이 저하되는 경향이 있다.\\
        
        \subsubsection{Computational Cost}
            표 \ref{tab:compute_cost}에 각 필터의 주요 연산량을 비교하였다.\\
 
            \begin{table}[H]
                \centering
                \caption{필터별 연산량 비교}
                \label{tab:compute_cost}
                \begin{tabular}{lcccc}
                    \toprule
                    필터 & 곱셈 & 덧셈 & 행렬 역산 & 적합 환경 \\
                    \midrule
                    CF       & $O(1)$ & $O(1)$ & 불필요 & 초저전력 MCU \\
                    Mahony   & $O(1)$ & $O(1)$ & 불필요 & 임베디드 \\
                    Madgwick & $O(1)$ & $O(1)$ & 불필요 & 임베디드 \\
                    EKF      & $O(n^2)$ & $O(n^2)$ & $3\times 3$ & 고성능 MCU / SBC \\
                    \bottomrule
                \end{tabular}
            \end{table}

            본 연구의 EKF는 상태 벡터 차원이 $n = 7$이므로, $7\times 7$ 행렬 연산과 $3\times 3$ 행렬 역산이 매 스텝 필요하다. $f_s = 100\,\text{Hz}$ 기준으로 Arduino Uno(16 MHz)에서는 실시간 실행이 어려울 수 있으며, Arduino Due 또는 Raspberry Pi 수준의 플랫폼이 적합하다.\\

        \subsubsection{Parameter Sensitivity}
            각 필터의 튜닝 파라미터에 대한 민감도를 평가한다.\\
    
            CF의 $\alpha$는 $[0.95, 0.99]$ 범위에서 성능 변화가 크지 않아 튜닝이 용이하다. Mahony의 $k_p$는 수렴 속도에 민감하며, 과도하게 크면 진동이 발생한다. EKF의 $\mathbf{Q}/\mathbf{R}$ 비율은 Allan Variance 결과로 합리적인 초기값을 설정할 수 있어 튜닝 부담이 상대적으로 낮다.\\

    \subsection{Limitations}
        \subsubsection{Sensor Noise and Drift}
            본 연구에서 사용한 MPU-9250은 소비자용 MEMS IMU로, 산업용 또는 항법용 IMU에 비해 ARW 및 BI 값이 크다. 특히 자이로 바이어스 드리프트는 온도 변화에 민감하여, 실내 환경에서도 수 분 내에 수 도($^\circ$)의 Yaw 오차가 발생할 수 있다.\\

        \subsubsection{Magnetic Disturbance}
            지자기계는 실내 환경에서의 자기 왜곡에 취약하다. Hard/Soft Iron 보정을 수행하더라도, 주변 전자 장비의 동작 상태에 따라 자기장이 변화하는 경우 보정 효과가 제한된다. 이러한 환경에서는 지자기계 Update를 선택적으로 비활성화하거나, 자기 왜곡 감지 알고리즘을 추가하는 것이 필요하다.\\

        \subsubsection{Dynamic Acceleration Effect}
            가속도계 기반 roll/pitch 추정은 정적 가속도 가정에 기반하므로, 선형 가속도가 큰 동적 환경에서는 오차가 증가한다. 본 연구에서는 임계값 기반 적응 조건(식 \eqref{eq:ekf_accel_condition})을 적용하였으나, 임계값 설정이 어렵고 빠른 가속 변화에 반응이 늦는 한계가 있다.\\

    \subsection{Future Work}
        \subsubsection{IMU Preintegration}
            현재 EKF는 각 타임스텝마다 독립적으로 센서 데이터를 처리한다. IMU Preintegration 기법을 적용하면 고주파 IMU 데이터를 키프레임(keyframe) 사이에서 효율적으로 적분하여 SLAM(Simultaneous Localization and Mapping)이나 Visual-Inertial Odometry(VIO) 시스템과의 결합이 가능해진다.\\

        \subsubsection{Unscented Kalman Filter (UKF)}
            EKF는 1차 선형화에 기반하므로 비선형성이 강한 시스템에서 추정 오차가 발생할 수 있다. Unscented Kalman Filter(UKF)는 Sigma Point를 이용하여 비선형 변환을 2차 정확도로 근사하므로 이론적으로 더 정확한 추정이 가능하다. 특히 초기 수렴 과정에서의 성능 개선이 기대된다.\\

        \subsubsection{Deep Learning-based Approaches}
            최근 LSTM, Transformer 등 딥러닝 기반 IMU 자세 추정 방법이 주목받고 있다. TLIO, RONIN 등의 연구에서는 IMU 데이터로부터 관성 오도메트리(inertial odometry)를 학습하여 드리프트 없는 위치 추정을 달성하였다. 향후 딥러닝 기반 방법과 EKF를 결합한 하이브리드 접근이 유망하다.\\

    \subsection{Conclusion}
        본 연구에서는 MPU-9250 IMU를 이용한 AHRS를 설계하고 구현하였다. 주요 기여는 다음과 같다.\\

        첫째, IMU 센서의 동작 원리와 확률적 잡음 모델을 체계적으로 정리하고, Allan Variance 및 PSD 분석을 통해 잡음 파라미터를 정량화하였다.\\

        둘째, 가속도계 정적 캘리브레이션과 지자기계 Hard/Soft Iron 보정을 ellipsoid fitting 기법으로 수행하여 측정값의 정확도를 개선하였다.\\

        셋째, 쿼터니언 기반 EKF를 설계하여 짐벌 락 없이 전역적으로 안정적인 자세 추정을 구현하였으며, Allan Variance 결과를 공정 잡음 공분산 $\mathbf{Q}$에 직접 연계함으로써 체계적인 파라미터 설정 방법을 제시하였다.\\
        
        넷째, Complementary Filter, Mahony, Madgwick Filter와의 비교를 통해 각 알고리즘의 정확도, 계산 비용, 튜닝 용이성 측면에서의 트레이드오프를 분석하였다.\\

        시뮬레이션 및 실험 결과, EKF는 정확도와 바이어스 추정 능력에서 다른 필터에 비해 우수한 성능을 보였으나, 실내 자기 왜곡 환경에서는 Yaw 추정의 한계가 확인되었다. 향후 UKF 적용, IMU Preintegration, 딥러닝 기반 융합 등의 연구를 통해 성능을 더욱 개선할 수 있을 것으로 기대된다.\\